The surface representation in \eqref{eq:boundary-curvature} is useful for
theory but direct numerical integration becomes impractical as the observation
dimension $p$ grows.  We instead estimate its equivalent Dirac-delta
representation using observations from $P_0$.  This avoids both a spatial grid
and the incorrect substitution of the fitted CML mixture for the unknown
data-generating density.

For $j<l$, define the active-pair indicator
\begin{equation}
 \chi_{jl}(\mathbf{x};\boldsymbol{\theta})
 =\ind\!\left\{
 \min(q_j(\mathbf{x};\boldsymbol{\theta}),q_l(\mathbf{x};\boldsymbol{\theta}))
 >\max_{m\notin\{j,l\}}q_m(\mathbf{x};\boldsymbol{\theta})
 \right\},
 \label{eq:active-pair-indicator}
\end{equation}
with the indicator equal to one when $k=2$.  Let $K$ be a kernel and $h=h_n$
a bandwidth.  The proposed estimator is
\begin{equation}
 \widehat{\mathbf{B}}_n
 =\frac{1}{nh}
 \sum_{i=1}^n\sum_{1\leq j<l\leq k}
 K\!\left\{\frac{g_{jl}(\mathbf{X}_i;\widehat{\boldsymbol{\theta}}_n)}{h}\right\}
 \chi_{jl}(\mathbf{X}_i;\widehat{\boldsymbol{\theta}}_n)
 \nabla_{\boldsymbol{\theta}}g_{jl}(\mathbf{X}_i;\widehat{\boldsymbol{\theta}}_n)
 \nabla_{\boldsymbol{\theta}}g_{jl}(\mathbf{X}_i;\widehat{\boldsymbol{\theta}}_n)^\mathsf T.
 \label{eq:boundary-estimator}
\end{equation}
Its cost is linear in $n$ and quadratic in $k$, and it does not require a grid
in $\mathbb R^p$.

\begin{assumption}[Kernel and tube regularity]
\label{ass:kernel-tube}
\begin{enumerate}
 \item $K$ is bounded, Lipschitz, symmetric, compactly supported, and
 $\int K(v)\,dv=1$.
 \item $h\to0$, $nh\to\infty$, and
 $\|\widehat{\boldsymbol{\theta}}_n-\boldsymbol{\theta}^*\|/h=o_p(1)$.
 \item In neighbourhoods of active faces, the conditional surface moments of
 $\nabla_{\boldsymbol{\theta}}g_{jl}\nabla_{\boldsymbol{\theta}}g_{jl}^\mathsf T$ given $g_{jl}=v$ are continuous at
 $v=0$ and dominated by an integrable envelope.
 \item Contributions from $h$-neighbourhoods of multiple-tie junctions are
 $o(1)$ after the $h^{-1}$ normalization.
\end{enumerate}
\end{assumption}

For a root-$n$ consistent estimator, the plug-in condition in part (2) follows
from $\sqrt n\,h\to\infty$.  The familiar choice $h\asymp n^{-1/5}$ satisfies
the three rate restrictions, although data-driven bandwidth selection and
sensitivity analysis remain necessary in applications.

To describe precision, define the oracle estimator
$\widetilde{\mathbf{B}}_n$ by replacing
$\widehat{\boldsymbol{\theta}}_n$ with $\boldsymbol{\theta}^*$ in
\eqref{eq:boundary-estimator}.  For an active pair, let
\begin{equation}
 \boldsymbol{\Gamma}_{jl}(v)
 =\int_{\{\mathbf{x}:g_{jl}(\mathbf{x};\boldsymbol{\theta}^*)=v\}}
 \frac{p_0(\mathbf{x})\chi_{jl}(\mathbf{x};\boldsymbol{\theta}^*)
 \mathbf{W}_{jl}(\mathbf{x})}
 {\|\nabla_{\mathbf{x}}g_{jl}(\mathbf{x};\boldsymbol{\theta}^*)\|}
 \,d\mathcal H^{p-1}(\mathbf{x}),
 \label{eq:surface-moment}
\end{equation}
where
\begin{equation}
 \mathbf{W}_{jl}(\mathbf{x})
 =\nabla_{\boldsymbol{\theta}}g_{jl}(\mathbf{x};\boldsymbol{\theta}^*)
  \nabla_{\boldsymbol{\theta}}g_{jl}(\mathbf{x};\boldsymbol{\theta}^*)^\mathsf T.
 \label{eq:boundary-outer-product}
\end{equation}

\begin{theorem}[Bias, variance, and oracle limit distribution]
\label{thm:boundary-oracle-clt}
In addition to Assumption~\ref{ass:kernel-tube}, suppose each
$\boldsymbol{\Gamma}_{jl}$ is twice continuously differentiable at zero,
$K$ is a second-order kernel, and the corresponding fourth-order surface
moments are finite.  Let
$\mu_2(K)=\int v^2K(v)\,dv$ and $R(K)=\int K(v)^2\,dv$.  Then
\begin{align}
 E\widetilde{\mathbf{B}}_n
 &=\mathbf{B}+\frac{h^2\mu_2(K)}2
   \sum_{j<l}\boldsymbol{\Gamma}_{jl}''(0)+o(h^2),
 \label{eq:boundary-bias}\\
 \operatorname{Var}\{\operatorname{vec}(\widetilde{\mathbf{B}}_n)\}
 &=\frac1{nh}\boldsymbol{\Omega}_{\mathbf{B}}
   +o\{(nh)^{-1}\},
 \label{eq:boundary-variance}
\end{align}
where
\begin{equation}
\begin{aligned}
 \boldsymbol{\Omega}_{\mathbf{B}}
 &=R(K)\sum_{j<l}\int_{\mathcal F_{jl}^*}
 \frac{p_0(\mathbf{x})}
 {\|\nabla_{\mathbf{x}}g_{jl}(\mathbf{x};\boldsymbol{\theta}^*)\|}\\
 &\quad\times
 \operatorname{vec}\{\mathbf{W}_{jl}(\mathbf{x})\}
 \operatorname{vec}\{\mathbf{W}_{jl}(\mathbf{x})\}^\mathsf T
 \,d\mathcal H^{p-1}(\mathbf{x}).
\end{aligned}
 \label{eq:boundary-asymptotic-variance}
\end{equation}
If $nh^5\to0$, then
\begin{equation}
 \sqrt{nh}\,\operatorname{vec}
 (\widetilde{\mathbf{B}}_n-\mathbf{B})
 \rightsquigarrow N_{d^2}(\mathbf{0},\boldsymbol{\Omega}_{\mathbf{B}}).
 \label{eq:boundary-oracle-clt}
\end{equation}
Without undersmoothing, the same conclusion holds after centering by the
leading bias in \eqref{eq:boundary-bias}.
\end{theorem}

\begin{proof}
The coarea formula writes the expectation as
$\sum_{j<l}\int K(t)\boldsymbol{\Gamma}_{jl}(ht)\,dt$.
A second-order Taylor expansion gives \eqref{eq:boundary-bias}.  Applying the
same representation to the squared vectorized summand gives
\eqref{eq:boundary-variance}; cross-pair terms are lower order because
first-order mass at multiple-face junctions is excluded.  A triangular-array
Lindeberg argument yields the centered normal limit.  Finally,
$\sqrt{nh}h^2\to0$ is equivalent to $nh^5\to0$ and removes the leading bias.
\end{proof}

\subsection{Finite-sample curvature stabilization}

Kernel variability can make
$\widehat{\mathbf{A}}_n=\widehat{\mathbf{I}}_{c,n}-
\widehat{\mathbf B}_n$ nonpositive in finite samples even when
$\mathbf{A}$ is positive definite.  To make this failure visible while still
providing usable inference, define
\begin{align}
 \widehat\lambda_n
 &=\lambda_{\max}\!\left(
 \widehat{\mathbf I}_{c,n}^{-1/2}
 \widehat{\mathbf B}_n
 \widehat{\mathbf I}_{c,n}^{-1/2}\right),                 \label{eq:relative-boundary-eigenvalue}\\
 \widehat\gamma_n
 &=\min\left\{1,\frac{1-\varepsilon_n}{\widehat\lambda_n}\right\},
 \qquad
 \widehat{\mathbf B}_n^{\mathrm{stab}}
 =\widehat\gamma_n\widehat{\mathbf B}_n,                 \label{eq:boundary-stabilization}
\end{align}
where $\varepsilon_n\downarrow0$.  We use
$\varepsilon_n=n^{-1/4}$ in the simulations.

\begin{proposition}[Asymptotic inactivity of stabilization]
\label{prop:stabilization-inactive}
Suppose $\mathbf{I}_c$ and
$\mathbf{A}=\mathbf{I}_c-\mathbf{B}$ are positive definite and the conditions of
Theorem~\ref{thm:sandwich-consistency} hold.  If
$\varepsilon_n\to0$, then
\begin{equation}
 P(\widehat\gamma_n=1)\to1,
 \qquad
 \widehat{\mathbf{B}}_n^{\mathrm{stab}}\to_p\mathbf{B}.
 \label{eq:stabilization-inactive}
\end{equation}
Consequently, replacing $\widehat{\mathbf B}_n$ by its stabilized version does
not change the first-order limiting distribution.
\end{proposition}

\begin{proof}
Positive definiteness of $\mathbf{A}$ implies
\[
 \lambda_{\max}(\mathbf{I}_c^{-1/2}\mathbf{B}
                         \mathbf{I}_c^{-1/2})<1.
\]
Consistency of $\widehat{\mathbf I}_{c,n}$ and
$\widehat{\mathbf B}_n$, together with continuity of inverse square roots and
eigenvalues on the positive-definite cone, gives
$\widehat\lambda_n\to_p\lambda_{\max}<1$.  Since
$\varepsilon_n\to0$, the inequality
$\widehat\lambda_n\leq1-\varepsilon_n$ holds with probability tending to one.
Hence $\widehat\gamma_n=1$ with probability tending to one, proving both
claims.
\end{proof}

\begin{proposition}[Plug-in error rate]
\label{prop:boundary-plugin-rate}
Under the conditions of Theorem~\ref{thm:boundary-oracle-clt}, suppose the
probability contribution of $h$-neighbourhoods of multiple-tie junctions,
after $h^{-1}$ normalization, is $O(\zeta_h)$ with $\zeta_h\to0$.  Then
\begin{equation}
 \|\widehat{\mathbf{B}}_n-\mathbf{B}\|_F
 =O_p\!\left(
 h^2+(nh)^{-1/2}
 +\frac{\|\widehat{\boldsymbol{\theta}}_n
             -\boldsymbol{\theta}^*\|}{h}
 +\zeta_h
 \right).
 \label{eq:boundary-plugin-rate}
\end{equation}
For root-$n$ consistent $\widehat{\boldsymbol{\theta}}_n$ and
$h\asymp n^{-1/5}$, all displayed terms converge to zero.
\end{proposition}

\begin{proof}
The first two terms follow from
Theorem~\ref{thm:boundary-oracle-clt}.  Lipschitz continuity of $K$, the
$h^{-1}$ scaling in its argument, and local derivative envelopes bound the
parameter-substitution error by the third term.  Classification of an active
pair can differ only near a multiple-tie junction or within a tube whose width
is proportional to the parameter error; these contributions are controlled by
$\zeta_h$ and the same substitution term.
\end{proof}

\begin{theorem}[Consistency of the boundary-curvature estimator]
\label{thm:boundary-estimator-consistency}
Under Assumptions~\ref{ass:boundary-geometry},
\ref{ass:local-stochastic}, and \ref{ass:kernel-tube},
\begin{equation}
 \widehat{\mathbf{B}}_n\xrightarrow{p}\mathbf{B}.
 \label{eq:boundary-estimator-consistency}
\end{equation}
\end{theorem}

\begin{proof}
First fix $\boldsymbol{\theta}=\boldsymbol{\theta}^*$.  The coarea formula gives, for each active pair,
\begin{multline*}
 \frac1hP_0\left[
 K\{g_{jl}(\mathbf{X};\boldsymbol{\theta}^*)/h\}\chi_{jl}(\mathbf{X};\boldsymbol{\theta}^*)
 \nabla_{\boldsymbol{\theta}}g_{jl}(\mathbf{X};\boldsymbol{\theta}^*)\nabla_{\boldsymbol{\theta}}g_{jl}(\mathbf{X};\boldsymbol{\theta}^*)^\mathsf T
 \right]\\
 =\int K(v/h)h^{-1}
 \int_{\{g_{jl}=v\}}
 \frac{p_0(\mathbf{x})\chi_{jl}(\mathbf{x};\boldsymbol{\theta}^*)
 \nabla_{\boldsymbol{\theta}}g_{jl}(\mathbf{x};\boldsymbol{\theta}^*)\nabla_{\boldsymbol{\theta}}g_{jl}(\mathbf{x};\boldsymbol{\theta}^*)^\mathsf T}
 {\|\nabla_{\mathbf{x}}g_{jl}(\mathbf{x};\boldsymbol{\theta}^*)\|}
 \,d\mathcal H^{p-1}(\mathbf{x})\,dv.
\end{multline*}
Continuity of the inner surface moment and $\int K=1$ make this converge to
the $(j,l)$ term of \eqref{eq:boundary-curvature}.  The variance is
$O\{(nh)^{-1}\}$ under the envelope conditions, so the centered empirical
term vanishes.  The condition
$\|\widehat{\boldsymbol{\theta}}_n-\boldsymbol{\theta}^*\|/h=o_p(1)$, Lipschitz regularity, and negligible
multiple-tie neighbourhoods make substitution of $\widehat{\boldsymbol{\theta}}_n$ and the
estimated active indicator asymptotically negligible.  Summing over the
finitely many pairs proves the result.
\end{proof}

The choice $h\asymp n^{-1/5}$ is used here for consistency and conventional
MSE-type bias--variance behaviour, not for an undersmoothed centred CLT.

Define the feasible fixed-classification and score matrices
\begin{align}
 \widehat{\mathbf{I}}_{c,n}
 &=-\frac1n\sum_{i=1}^n\sum_{j=1}^k
 \ind\{a_{\widehat{\boldsymbol{\theta}}_n}(\mathbf{X}_i)=j\}
 \nabla_{\boldsymbol{\theta}}^2q_j(\mathbf{X}_i;\widehat{\boldsymbol{\theta}}_n),                         \label{eq:ic-estimator}\\
 \widehat{\mathbf{J}}_n
 &=\frac1n\sum_{i=1}^n
 \widehat{\boldsymbol{\psi}}_i\widehat{\boldsymbol{\psi}}_i^\mathsf T,
 \qquad
 \widehat{\boldsymbol{\psi}}_i=\boldsymbol{\psi}(\mathbf{X}_i;\widehat{\boldsymbol{\theta}}_n),              \label{eq:j-estimator}\\
 \widehat{\mathbf{A}}_n
 &=\widehat{\mathbf{I}}_{c,n}-\widehat{\mathbf{B}}_n,               \label{eq:a-estimator}\\
 \widehat{\mathbf{V}}_n
 &=\widehat{\mathbf{A}}_n^{-1}\widehat{\mathbf{J}}_n\widehat{\mathbf{A}}_n^{-\mathsf T}.
 \label{eq:v-estimator}
\end{align}

\begin{theorem}[Consistency of feasible sandwich inference]
\label{thm:sandwich-consistency}
Suppose the assumptions of
Theorem~\ref{thm:boundary-estimator-consistency} hold, the interior Hessian and
score classes satisfy uniform laws of large numbers near $\boldsymbol{\theta}^*$, and $\mathbf{A}$
is nonsingular.  Then
\begin{equation}
 \widehat{\mathbf{I}}_{c,n}\to_p \mathbf{I}_c,\qquad
 \widehat{\mathbf{J}}_n\to_p \mathbf{J},\qquad
 \widehat{\mathbf{A}}_n\to_p \mathbf{A},\qquad
 \widehat{\mathbf{V}}_n\to_p \mathbf{V}.
 \label{eq:sandwich-consistency}
\end{equation}
Consequently, $\widehat{\mathbf{V}}_n/n$ consistently estimates the first-order
covariance matrix of $\widehat{\boldsymbol{\theta}}_n$.
\end{theorem}

\begin{proof}
Consistency of $\widehat{\boldsymbol{\theta}}_n$, the local uniform laws, and the fact that
ties have probability zero give consistency of $\widehat{\mathbf{I}}_{c,n}$ and
$\widehat{\mathbf{J}}_n$.  Theorem~\ref{thm:boundary-estimator-consistency} gives the
remaining component of $\widehat{\mathbf{A}}_n$.  Matrix inversion is continuous on the
set of nonsingular matrices, so the final assertion follows from the
continuous mapping theorem.
\end{proof}

\begin{corollary}[First-order equivalence after stabilization]
\label{cor:stabilized-equivalence}
Under Proposition~\ref{prop:stabilization-inactive} and
Theorem~\ref{thm:sandwich-consistency}, let
$\widehat{\mathbf V}_n^{\mathrm{stab}}$ be obtained from
$\widehat{\mathbf B}_n^{\mathrm{stab}}$.  Then
$P(\widehat{\mathbf V}_n^{\mathrm{stab}}=\widehat{\mathbf V}_n)\to1$ and
$\widehat{\mathbf V}_n^{\mathrm{stab}}\to_p\mathbf V$.  Hence stabilized and
unstabilized Wald statistics have the same first-order limit whenever
stabilization is asymptotically inactive.
\end{corollary}

\begin{corollary}[Wald inference]
\label{cor:wald-inference}
Let $\mathbf{R}$ be a fixed $r\times d$ matrix of rank $r$ and suppose
$\mathbf{R}\mathbf{V}\mathbf{R}^\mathsf T$ is nonsingular.  Under
$H_0:\mathbf{R}\boldsymbol{\theta}^*=\mathbf{c}$,
\begin{equation}
 n(\mathbf{R}\widehat{\boldsymbol{\theta}}_n-\mathbf{c})^\mathsf T
 (\mathbf{R}\widehat{\mathbf{V}}_n\mathbf{R}^\mathsf T)^{-1}
 (\mathbf{R}\widehat{\boldsymbol{\theta}}_n-\mathbf{c})
 \rightsquigarrow\chi_r^2.
 \label{eq:wald-statistic}
\end{equation}
In particular, componentwise intervals based on
$\{\widehat{\mathbf{V}}_{n,ss}/n\}^{1/2}$ have asymptotic coverage equal to their
nominal level.
\end{corollary}

\begin{corollary}[Smooth transformations]
\label{cor:delta-method}
Let $\boldsymbol{\tau}:\mathbb R^d\to\mathbb R^s$ be continuously
differentiable at $\boldsymbol{\theta}^*$ with Jacobian matrix
$\mathbf{D}_{\tau}$.  Then
\begin{equation}
 \sqrt n\{\boldsymbol{\tau}(\widehat{\boldsymbol{\theta}}_n)
 -\boldsymbol{\tau}(\boldsymbol{\theta}^*)\}
 \rightsquigarrow
 N_s(\mathbf{0},\mathbf{D}_{\tau}\mathbf{V}
                    \mathbf{D}_{\tau}^\mathsf T),
 \label{eq:delta-method}
\end{equation}
and the plug-in covariance obtained by replacing
$\mathbf{D}_{\tau}$ and $\mathbf{V}$ by consistent estimators is consistent.
This result covers scientifically interpretable transformations of mixture
weights, locations, covariance parameters, and component contrasts.
\end{corollary}
